\section*{Chapter \ref{ch:in:choosing} --- Choosing}

\begin{solution}{exo:in:choosing:1}
The technology engineer's offer, by the bank quant's: more expected pay, less spread, a smaller share taken by tax and
housing, no on-call, and equal learning and security.
\end{solution}

\begin{solution}{exo:in:choosing:2}
Expected pay 0.30, pay risk 0.15, place 0.10, on-call 0.15, learning 0.20, security 0.10.
\end{solution}

\begin{solution}{exo:in:choosing:3}
The market maker's offer on both: a certainty equivalent of \$1\,365\,186 and expected pay of \$1\,949\,485 over five
years.
\end{solution}

\begin{solution}{exo:in:choosing:4}
It has the widest spread of pay and the only chance of a closed job, the worst scores on two attributes, and it trails the
market maker on expected pay and on place; only draws that weight learning and no-on-call heavily and pay risk and
security lightly put it first.
\end{solution}

\begin{solution}{exo:in:choosing:5}
It treats every weighting as equally likely, including extreme ones no real chooser holds, and ignores what the chooser
knows about her own priorities; the acceptabilities describe the offers' robustness, not her preference.
\end{solution}

\begin{solution}{exo:in:choosing:6}
Still the market maker's offer (0.894, against 0.706 for the platform's): it ties for the best learning score.
\end{solution}

\begin{solution}{exo:in:choosing:7}
Rescaled over three offers, the first-rank acceptabilities become 81.4\% (market maker), 13.5\% (bank quant) and 5.1\%
(platform analyst): dropping a dominated offer changes the scales and so, a little, the shares.
\end{solution}

\begin{solution}{exo:in:choosing:8}
The weights are drawn uniformly by assumption, not from a survey of people; the index measures how much of the weight
space favours an offer, which says nothing about how many people hold which weights. And the attribute values are
illustrative.
\end{solution}

\begin{solution}{pb:in:choosing:1}
\begin{enumerate}
\item Pay and its spread (13, 14), place (15, 27), hours (26), the work and what it leads to (16--25, 28), security (22).
\item As in the chapter's definition.
\item \$1\,186\,354, \$1\,365\,186, \$903\,603 and \$845\,964.
\item So that the weight on risk is explicit and can be varied.
\item Learning, and the weights themselves.
\item As in the chapter's table.
\item As in the definition; the technology engineer's offer.
\item 30, 15, 10, 15, 20, 10; values 0.852 (market maker), 0.589 (platform), 0.584 (bank), 0.202 (technology).
\item As in the definition; scale, screen, weight, sample weights on the simplex and count ranks.
\item 79.5\%, 15.9\%, 4.6\% and 0\%.
\item By 0.452, from 0.15 to 0.602.
\item It is the choice only for someone who weights pay risk above all the other attributes together.
\item Evidence about what each job teaches: its \omterm{def:in:quant-researcher:feedback}{feedback speed}, its training, where its people go next.
\item It would raise the platform's expected pay and spread together.
\item Temperament, the people, and the chooser's circumstances.
\item Market maker 79.5\%, bank 15.9\%, platform 4.6\%, technology 0\%; the weight on pay risk must rise by 0.452.
\item In a decision journal: attributes, weights, expectations and what would change her mind.
\item When a belief changes (pay, learning, security), and at a date fixed in advance, not after a bad month.
\item A result she could not have foreseen is bad luck; a result she could have foreseen and ignored is a bad choice; the
journal tells them apart.
\item On the chapter's numbers the market maker's offer wins under most weightings, and only a very strong aversion to
pay risk favours the bank's. Record why, and revisit when a belief, not an outcome, changes.
\end{enumerate}
\end{solution}

\begin{solution}{iq:in:choosing:1}
Name the attributes that decide it for you (the work, what you will learn, the people) and show you know the job's
realities; do not run down the other offers.
\iqlookfor{specific reasons tied to the job.}
\end{solution}

\begin{solution}{iq:in:choosing:2}
By risk and timing: certainty equivalents under a stated risk aversion, deferral and \omterm{def:in:how-pay-works:rsu}{forfeiture} (what leaving costs), and
tax timing.
\iqlookfor{structure matters even at equal expected value.}
\end{solution}

\begin{solution}{iq:in:choosing:3}
Compare each pair: offer $i$ is dominated if some $j$ is at least as good everywhere and better somewhere; $O(n^2m)$ for
$n$ offers and $m$ attributes, fine for small $n$; a sort-based skyline algorithm does better for large $n$.
\iqlookfor{the pairwise definition and its cost.}
\end{solution}

\begin{solution}{iq:in:choosing:4}
Draw independent exponential variables and normalise them (a Dirichlet with unit parameters), or sort uniform draws and
take the gaps; normalising uniform draws is not uniform on the simplex.
\iqlookfor{the Dirichlet construction and the common mistake.}
\end{solution}

\begin{solution}{iq:in:choosing:5}
Compare the outcome with what was known and expected when choosing (the journal): if the bad result was a foreseeable
risk you accepted knowingly, it was bad luck; if you ignored information you had, it was a bad decision.
\iqlookfor{avoiding outcome bias.}
\end{solution}

\begin{solution}{iq:in:choosing:6}
The additive model double counts what the two attributes share, giving it too much weight; merge them, reduce the weights
of both, or keep only the one that matters to you.
\iqlookfor{double counting in additive models.}
\end{solution}
